\section{Inference in FOL}

\todo{Substitutions are used to lift inference to propositional logics:
    \begin{itemize}
        \item Represent the formulas to be infered with the same computation network of atomic formulas.
        \item Represent the formulas with an entangled activation network (whereas before they were disentangled, since refering to different atomic sentences).
    \end{itemize}
}



\subsection{Manipulation of contraction equations}

Universally quantified FOL formulas are equations of FOL \ComputationActivationNetworks{} to the trivial tensor.


We understand inference as manipulations in the activation network.
\begin{itemize}
    \item {\bf Substituion}: Substitution of variables to represent multiple sentences by same activation network.
    The task of unification is to find a substitution, such that two computation networks can be merged with same term variables.
    \item {\bf Propositional inference}: Manipulations in activation network, as for constraint networks
    \item {\bf Restrictions to parts}: Any sub-network of the activation network in a contraction equation poses an implied contraction equation.
\end{itemize}

\subsubsection{Substitution}

Substitution is the contraction with one-hot encodings to the term variables.
It is easy to see, that when a contraction equation holds, it also holds for any substitution of the term variables.
Note that substitution changes the computation network, but not the activation network.

\subsubsection{Propositional inference}

We understand propositional inference as the simplication of the activation network.
Whenever a propositional formula is infered by the activation network, we can add it to the activation network as a new decorated edge.

\subsubsection{Extraction of entailed formulas}

We can extract entailed formulas from a contraction equation by restricting the activation network to a sub-network, as we show as the next theorem.
Therefore any result of the propositional inference, in the shape of a Boolean tensor added to the activation network, can be used to extract entailed formulas.

\begin{theorem}
    Let $\extnet$ be an activation network of boolean tensors on a hypergraph $\graph$ and let $\secedges\subset\edges$.
    Let $\folpredicateof{[\catorder]}$ be an arbitrary set of first order predicates with the term variables
    If
    \begin{align*}
        \contractionof{\extnet\cup
            \{\bencodingofat{\catenumerator}{\catvariableof{\catenumerator},\indvariableof{[n]}} \wcols \catenumeratorin\}
        }{\indvariableof{[n]}}
        = \onesat{\catvariableof{[n]}}
    \end{align*}
    then also
    \begin{align*}
        \contractionof{
            \{\hypercoreofat{\edge}{\catvariableof{\edge}}\wcols\edge\in\secedges\}
            \cup
            \{\bencodingofat{\catenumerator}{\catvariableof{\catenumerator},\indvariableof{[n]}} \wcols \catenumeratorin\}
        }{\indvariableof{[n]}}
        = \onesat{\catvariableof{[n]}}
    \end{align*}
\end{theorem}
\begin{proof}
    Notice, that the contraction with the smaller set of edges $\secedges\subset\edges$ is a Boolean tensor.
    By the monotonicity of the contraction, we have that the support of the contraction is smaller (in the partial order of tensors) to that of all edges.
    Since by assumption the latter is the ones tensors, the former must also be the ones tensor.
\end{proof}




\subsection{Unification}

Unification of atomic sentences: Find a substitution to the variables, such that both sentences are represented by the same computation network.

Compressing of complex sentences:
\begin{itemize}
    \item Find a unification of the atomic sentences in both complex sentences
    \item Represent the substituted complex sentences by a constraction of the activation networks
\end{itemize}

\subsection{Modus ponens}

\begin{example}[Lifted modus ponens]\label{exa:liftedModusPonens}
    For example we have
    \begin{itemize}
        \item John knows everyone: Knows(John,z)
        \item People who know each other are friends: Knows(x,y) $\Rightarrow$ Friends(x,y)
    \end{itemize}
    We illustrate the derivation of the entailed formula Friends(John,z) (i.e. John is a friend of everyone) by modus ponens.
    The formulas are represented by the contraction equation

    \stantikzpic{
        \drawsmalltensor{5}{4}{$\tbasis$}
        \drawtdwire{5}{3}{$\catvariableof{Knows}$}
        \drawtensorblock{5}{0}{2}{$\bencodingof{Knows}$}
        \drawtdwire{3.5}{-1}{$\indvariableof{0}$}
        \drawtensorblock{3.5}{-4}{1.5}{$\onehotmapof{John}$}
        \drawtdwire{6.5}{-1}{$\indvariableof{1}$}

        \drawtextnode{10}{0}{$=$}

        \drawtensorblock{14}{0}{1}{$\ones$}
        \drawtdwire{14}{-1}{$\indvariableof{1}$}

        \drawtextnode{19}{0}{and}

        \begin{scope}[shift={(25,0)}]

            \drawtensorblock{2.5}{4}{3}{$\impbincon$}
            \drawtdwire{0}{3}{$\catvariableof{Knows}$}
            \drawtensorblock{0}{0}{2}{$\bencodingof{Knows}$}
            \drawtdwire{5}{3}{$\catvariableof{Friends}$}
            \drawtensorblock{5}{0}{2}{$\bencodingof{Friends}$}

            \draw (1,-3) to[bend right=-20] (-1,-1);
            \draw (1,-3) to[bend right=20] (4,-1);
            \drawvariabledot{1}{-3}
            \drawtdwire{1}{-3}{$\indvariableof{2}$}
            \draw (4,-3) to[bend right=-20] (1,-1);
            \draw (4,-3) to[bend right=20] (6,-1);
            \drawvariabledot{4}{-3}
            \drawtdwire{4}{-3}{$\indvariableof{3}$}

            \drawtextnode{10}{0}{$=$}

            \drawtensorblock{15}{0}{2}{$\ones$}
            \drawtdwire{13.5}{-1}{$\indvariableof{2}$}
            \drawtdwire{16.5}{-1}{$\indvariableof{3}$}
        \end{scope}
    }

    We combine them by the tensor product of both sides as

    \stantikzpic{

        \drawsmalltensor{5}{4}{$\tbasis$}
        \drawtdwire{5}{3}{$\catvariableof{Knows,0}$}
        \drawtensorblock{5}{0}{2}{$\bencodingof{Knows}$}
        \drawtdwire{3.5}{-1}{$\indvariableof{0}$}
        \drawtensorblock{3.5}{-4}{1.5}{$\onehotmapof{John}$}
        \drawtdwire{6.5}{-1}{$\indvariableof{1}$}

        \begin{scope}[shift={(10,0)}]

            \drawtensorblock{2.5}{4}{3}{$\impbincon$}
            \drawtdwire{0}{3}{$\catvariableof{Knows,1}$}
            \drawtensorblock{0}{0}{2}{$\bencodingof{Knows}$}
            \drawtdwire{5}{3}{$\catvariableof{Friends}$}
            \drawtensorblock{5}{0}{2}{$\bencodingof{Friends}$}

            \draw (1,-3) to[bend right=-20] (-1,-1);
            \draw (1,-3) to[bend right=20] (4,-1);
            \drawvariabledot{1}{-3}
            \drawtdwire{1}{-3}{$\indvariableof{2}$}
            \draw (4,-3) to[bend right=-20] (1,-1);
            \draw (4,-3) to[bend right=20] (6,-1);
            \drawvariabledot{4}{-3}
            \drawtdwire{4}{-3}{$\indvariableof{3}$}


        \end{scope}


        \drawtextnode{19}{0}{$=$}

        \drawtensorblock{23}{0}{2.5}{$\ones$}
        \drawtdwire{21}{-1}{$\indvariableof{1}$}
        \drawtdwire{23}{-1}{$\indvariableof{2}$}
        \drawtdwire{25}{-1}{$\indvariableof{3}$}
    }

    The necessary substitution for the unification of both is $\indvariableof{2}\mapsto$ John and $\indvariableof{1},\indvariableof{3}\mapsto\indvariableof{4}$.
    This amounts to contraction of both sides with the tensor
    \stantikzpic{
        \drawtdwire{-1}{3}{$\indvariableof{1}$}
        \drawtdwire{1}{3}{$\indvariableof{3}$}
        \drawtensorblock{0}{0}{1.5}{$\identity$}
        \drawtdwire{0}{-1}{$\indvariableof{4}$}

        \drawtdwire{3.5}{3}{$\indvariableof{2}$}
        \drawtensorblock{3.5}{0}{1.5}{$\onehotmapof{John}$}
    }

    We contract this tensor to both sides of the combined contraction equation and get
    \stantikzpic{
        \drawsmalltensor{-5}{4}{$\tbasis$}
        \drawtdwire{-5}{3}{$\catvariableof{Knows,0}$}
        \drawtensorblock{-5}{0}{2}{$\bencodingof{Knows}$}

        \drawtensorblock{2.5}{4}{3}{$\impbincon$}
        \drawtdwire{0}{3}{$\catvariableof{Knows,1}$}
        \drawtensorblock{0}{0}{2}{$\bencodingof{Knows}$}
        \drawtdwire{5}{3}{$\catvariableof{Friends}$}
        \drawtensorblock{5}{0}{2}{$\bencodingof{Friends}$}


        \draw (1,-3) to[bend right=-10] (-6,-1);
        \draw (1,-3) to[bend right=-20] (-1,-1);
        \draw (1,-3) to[bend right=20] (4,-1);
        \drawvariabledot{1}{-3}
        \drawtdwire{1}{-3}{$\indvariableof{2}$}
        \drawtensorblock{1}{-6}{1.5}{$\onehotmapof{John}$}


        \draw (4,-3) to[bend right=-10] (-4,-1);
        \draw (4,-3) to[bend right=-20] (1,-1);
        \draw (4,-3) to[bend right=20] (6,-1);
        \drawvariabledot{4}{-3}
        \drawtdwire{4}{-3}{$\indvariableof{4}$}

        \drawtextnode{9}{0}{$=$}

        \drawtensorblock{13}{0}{1}{$\ones$}
        \drawtdwire{13}{-1}{$\indvariableof{4}$}
    }

    Now we notice that the computation network of the variables $\catvariableof{Knows,0}$ and $\catvariableof{Knows,1}$ are equivalent (what was the goal of the unification).
    We simplify the equation to
    \stantikzpic{
        \drawsmalltensor{0}{5}{$\tbasis$}
        \drawtensorblock{4}{5}{2}{$\impbincon$}

        \draw (0,3) to [bend right=20] (3,4);
        \draw (0,3) -- (0,4);

        \drawvariabledot{0}{3}
        \drawtdwire{0}{3}{$\catvariableof{Knows}$}
        \drawtensorblock{0}{0}{2}{$\bencodingof{Knows}$}

        \draw (5,3) -- (5,4);
        \drawvariabledot{5}{3}
        \drawtdwire{5}{3}{$\catvariableof{Friends}$}
        \drawtensorblock{5}{0}{2}{$\bencodingof{Friends}$}

        \draw (1,-3) to[bend right=-20] (-1,-1);
        \draw (1,-3) to[bend right=20] (4,-1);
        \drawvariabledot{1}{-3}
        \drawtdwire{1}{-3}{$\indvariableof{2}$}
        \drawtensorblock{1}{-6}{1.5}{$\onehotmapof{John}$}

        \draw (4,-3) to[bend right=-20] (1,-1);
        \draw (4,-3) to[bend right=20] (6,-1);
        \drawvariabledot{4}{-3}
        \drawtdwire{4}{-3}{$\indvariableof{4}$}

        \drawtextnode{9}{0}{$=$}

        \begin{scope}[shift={(13,0)}]

            \drawsmalltensor{0}{5}{$\tbasis$}
            \drawsmalltensor{5}{5}{$\tbasis$}

            \draw (0,3) -- (0,4);

            \drawvariabledot{0}{3}
            \drawtdwire{0}{3}{$\catvariableof{Knows}$}
            \drawtensorblock{0}{0}{2}{$\bencodingof{Knows}$}

            \draw (5,3) -- (5,4);
            \drawvariabledot{5}{3}
            \drawtdwire{5}{3}{$\catvariableof{Friends}$}
            \drawtensorblock{5}{0}{2}{$\bencodingof{Friends}$}

            \draw (1,-3) to[bend right=-20] (-1,-1);
            \draw (1,-3) to[bend right=20] (4,-1);
            \drawvariabledot{1}{-3}
            \drawtdwire{1}{-3}{$\indvariableof{2}$}
            \drawtensorblock{1}{-6}{1.5}{$\onehotmapof{John}$}

            \draw (4,-3) to[bend right=-20] (1,-1);
            \draw (4,-3) to[bend right=20] (6,-1);
            \drawvariabledot{4}{-3}
            \drawtdwire{4}{-3}{$\indvariableof{4}$}

            \drawtextnode{9}{0}{$=$}

            \drawtensorblock{13}{0}{1}{$\ones$}
            \drawtdwire{13}{-1}{$\indvariableof{4}$}
        \end{scope}

    }

    Here we applied modus ponens on the activation network in the first equation.
    Restricting the activation network to the variable $\catvariableof{Friends}$ is then the derived formula Friends(John,z) (i.e. John is a friend of everyone):

    \stantikzpic{
        \drawsmalltensor{5}{5}{$\tbasis$}

        \draw (5,3) -- (5,4);
        \drawvariabledot{5}{3}
        \drawtdwire{5}{3}{$\catvariableof{Friends}$}
        \drawtensorblock{5}{0}{2}{$\bencodingof{Friends}$}

        \drawtdwire{3.5}{-1}{$\indvariableof{2}$}
        \drawtensorblock{3.5}{-4}{1.5}{$\onehotmapof{John}$}

        \drawtdwire{6.5}{-1}{$\indvariableof{4}$}

        \drawtextnode{9}{0}{$=$}

        \drawtensorblock{13}{0}{1}{$\ones$}
        \drawtdwire{13}{-1}{$\indvariableof{4}$}

    }





% OLD, more direct sketch
%    Now we do modus ponens on the activation network and get Friends(John,z).
%    \stantikzpic{
%% From Knows(John,z)
%        \drawsquaretensor{0}{4}{$\tbasis$}
%        \draw (0,2) -- (0,3);
%
%        % From Knows(x,y) $\Rightarrow$ Friends(x,y)
%        \draw (3,3) rectangle (7,5);
%        \drawtextnode{5}{4}{\corelabelsize $\Rightarrow$};
%        \draw (0,2) to[bend right=20] (4,3);
%        \draw (6,2) -- (6,3);
%
%        \drawvariabledot{0}{2}
%        \drawvariabledot{6}{2}
%
%        %% SUBSTITUTION
%        \draw [->-] (0,1) -- (0,2) node[midway,right] {\colorlabelsize $\indvariableof{Knows}$};
%        \draw (-2,-1) rectangle (2,1);
%        \drawtextnode{0}{0}{\corelabelsize $\bencodingof{Knows}$};
%
%        \begin{scope}[shift={(6,0)}]
%            \draw [->-] (0,1) -- (0,2) node[midway,right] {\colorlabelsize $\indvariableof{Friends}$};
%            \draw (-2,-1) rectangle (2,1);
%            \drawtextnode{0}{0}{\corelabelsize $\bencodingof{Friends}$};
%        \end{scope}
%
%        \draw[->-] (1.5,-4) to[bend right=-40] (-1.5,-1);
%        \draw[->-] (1.5,-4) to[bend right=40] (4.5,-1);
%        \drawvariabledot{1.5}{-4}
%        \draw[->-] (1.5,-5) -- (1.5,-4) node[midway,right] {\colorlabelsize $\indvariableof{x}$};
%        \drawsquaretensor{1.5}{-6}{$\onehotmapof{John}$}
%
%        \draw[->-] (4.5,-4) to[bend right=-40] (1.5,-1);
%        \draw[->-] (4.5,-4) to[bend right=40] (7.5,-1);
%        \drawvariabledot{4.5}{-4}
%        \draw[->-] (4.5,-5) -- (4.5,-4) node[midway,right] {\colorlabelsize $\indvariableof{z}$};
%        \drawsquaretensor{4.5}{-6}{$\identity$}
%        \draw[->-] (4.5,-8.5) -- (4.5,-7) node[midway,right] {\colorlabelsize $\indvariableof{z}$};
%
%        \drawtextnode{13}{0}{${=}$}
%    %%% RIGHT SIDE
%        \begin{scope}[shift={(20,0)}]
%            % From Knows(John,z)
%            \drawsquaretensor{0}{4}{$\tbasis$}
%            \draw (0,2) -- (0,3);
%
%            % Modus ponens infered
%            \drawsquaretensor{6}{4}{$\tbasis$}
%            \draw (6,2) -- (6,3);
%
%            \drawvariabledot{0}{2}
%            \drawvariabledot{6}{2}
%
%            %% SUBSTITUTION
%            \draw [->-] (0,1) -- (0,2) node[midway,right] {\colorlabelsize $\indvariableof{Knows}$};
%            \draw (-2,-1) rectangle (2,1);
%            \drawtextnode{0}{0}{\corelabelsize $\bencodingof{Knows}$};
%
%            \begin{scope}[shift={(6,0)}]
%                \draw [->-] (0,1) -- (0,2) node[midway,right] {\colorlabelsize $\indvariableof{Friends}$};
%                \draw (-2,-1) rectangle (2,1);
%                \drawtextnode{0}{0}{\corelabelsize $\bencodingof{Friends}$};
%            \end{scope}
%
%            \draw[->-] (1.5,-4) to[bend right=-40] (-1.5,-1);
%            \draw[->-] (1.5,-4) to[bend right=40] (4.5,-1);
%            \drawvariabledot{1.5}{-4}
%            \draw[->-] (1.5,-5) -- (1.5,-4) node[midway,right] {\colorlabelsize $\indvariableof{x}$};
%            \drawsquaretensor{1.5}{-6}{$\onehotmapof{John}$}
%
%            \draw[->-] (4.5,-4) to[bend right=-40] (1.5,-1);
%            \draw[->-] (4.5,-4) to[bend right=40] (7.5,-1);
%            \drawvariabledot{4.5}{-4}
%            \draw[->-] (4.5,-5) -- (4.5,-4) node[midway,right] {\colorlabelsize $\indvariableof{z}$};
%            \drawsquaretensor{4.5}{-6}{$\identity$}
%            \draw[->-] (4.5,-8.5) -- (4.5,-7) node[midway,right] {\colorlabelsize $\indvariableof{z}$};
%        \end{scope}
%    }
\end{example}

\subsection{Contraction equations}

Quantified formulas understood are contraction equations.

Understanding them as hard constraints, one can infer parts of the substitution tensors, which holds for all supported worlds.

\begin{example}[Modus ponens on the substitution structure]
    \todo{Repetition?}
    We here examplify an application of unification to apply the modus ponens inference rule in the substitution structure.
    Let us consider the formula
    \begin{align*}
        \uniquantwrtof{\shortcatindices}{\formulaat{\indexedshortcatvariables}\impbincon\secformulaat{\indexedshortcatvariables}} \, ,
    \end{align*}
    which corresponds to the contraction equation
    \stantikzpic{
        \drawtensorblock{2.5}{4}{3}{$\impbincon$}
        \drawtdwire{0}{3}{$\catvariableof{\formula}$}
        \drawtensorblock{0}{0}{2}{$\bencodingof{\formula}$}
        \drawtdwire{5}{3}{$\catvariableof{\secformula}$}
        \drawtensorblock{5}{0}{2}{$\bencodingof{\secformula}$}

        \draw (1,-3) to[bend right=-20] (-1,-1);
        \draw (1,-3) to[bend right=20] (4,-1);
        \drawvariabledot{1}{-3}
        \drawtdwire{1}{-3}{$\indvariableof{0}$}
        \draw (4,-3) to[bend right=-20] (1,-1);
        \draw (4,-3) to[bend right=20] (6,-1);
        \drawvariabledot{4}{-3}
        \drawtdwire{4}{-3}{$\indvariableof{\indorder-1}$}
        \drawtextnode{3}{-4.5}{\colorlabelsize $\cdots$}

        \drawtextnode{10}{0}{$=$}

        \drawtensorblock{15}{0}{2}{$\ones$}
        \drawtdwire{13.5}{-1}{$\indvariableof{0}$}
        \drawtextnode{15.5}{-2.5}{\colorlabelsize $\cdots$}
        \drawtdwire{16.5}{-1}{$\indvariableof{\indorder-1}$}
    }
    When knowing for some $\shortcatindicesin$ that $\formulaat{\indexedshortcatvariables}=1$, this is the contraction equation
    we conclude that $\secformulaat{\indexedshortcatvariables}=1$.

    \stantikzpic{
        \drawtensorblock{2.5}{4}{3}{$\impbincon$}
        \drawtdwire{0}{3}{$\catvariableof{\formula}$}
        \drawtensorblock{0}{0}{2}{$\bencodingof{\formula}$}
        \drawtdwire{5}{3}{$\catvariableof{\secformula}$}
        \drawtensorblock{5}{0}{2}{$\bencodingof{\secformula}$}

        \draw (1,-3) to[bend right=-20] (-1,-1);
        \draw (1,-3) to[bend right=20] (4,-1);
        \drawvariabledot{1}{-3}
        \drawtdwire{1}{-3}{$\indvariableof{0}$}
        \draw (4,-3) to[bend right=-20] (1,-1);
        \draw (4,-3) to[bend right=20] (6,-1);
        \drawvariabledot{4}{-3}
        \drawtdwire{4}{-3}{$\indvariableof{\indorder-1}$}
        \drawtextnode{3}{-4.5}{\colorlabelsize $\cdots$}

        \drawsmalltensor{1}{-6}{$\onehotmapof{\catindexof{0}}$}
        \drawsmalltensor{4}{-6}{$\onehotmapof{\catindexof{\indorder\shortminus1}}$}

        \drawtextnode{10}{0}{$=$}

        \drawtensorblock{15}{0}{2}{$\ones$}
        \drawtdwire{13.5}{-1}{$\indvariableof{0}$}
        \drawtextnode{15.5}{-2.5}{\colorlabelsize $\cdots$}
        \drawtdwire{16.5}{-1}{$\indvariableof{\indorder-1}$}
        \drawsmalltensor{13.5}{-4}{$\onehotmapof{\catindexof{0}}$}
        \drawsmalltensor{16.5}{-4}{$\onehotmapof{\catindexof{\indorder\shortminus1}}$}


        \drawtextnode{20}{0}{$=\quad 1$}
    }

    We simplify the left hand side using basis calculus to

    \stantikzpic{
        \drawtensorblock{2.5}{4}{3}{$\impbincon$}
        \drawtdwire{0}{3}{$\catvariableof{\formula}$}
        \drawtensorblock{-1}{0}{2.5}{$\onehotmapof{\formulaat{\indexedshortcatvariables}}$}
        \drawtdwire{5}{3}{$\catvariableof{\secformula}$}
        \drawtensorblock{5}{0}{2}{$\bencodingof{\secformula}$}
        \drawtdwire{3.5}{-1}{$\indvariableof{0}$}
        \drawtextnode{5.5}{-2.25}{\colorlabelsize $\cdots$}
        \drawtdwire{6.5}{-1}{$\indvariableof{\indorder\shortminus1}$}
        \drawsmalltensor{3.5}{-4}{$\onehotmapof{\catindexof{0}}$}
        \drawsmalltensor{6.5}{-4}{$\onehotmapof{\catindexof{\indorder\shortminus1}}$}

        \drawtextnode{11}{0}{$= \quad 1$}
    }
    When known that $\formulaat{\indexedshortcatvariables}=1$, this further simplifies to

    \stantikzpic{
        \drawtensorblock{5}{4}{1}{$\tbasis$}
        \drawtdwire{5}{3}{$\catvariableof{\secformula}$}
        \drawtensorblock{5}{0}{2}{$\bencodingof{\secformula}$}
        \drawtdwire{3.5}{-1}{$\indvariableof{0}$}
        \drawtextnode{5.5}{-2.25}{\colorlabelsize $\cdots$}
        \drawtdwire{6.5}{-1}{$\indvariableof{\indorder\shortminus1}$}
        \drawsmalltensor{3.5}{-4}{$\onehotmapof{\catindexof{0}}$}
        \drawsmalltensor{6.5}{-4}{$\onehotmapof{\catindexof{\indorder\shortminus1}}$}

        \drawtextnode{11}{0}{$= \quad 1$}
    }
    which is equal to $\secformulaat{\indexedshortcatvariables}=1$.
\end{example}

More sophisticated reasoning schemes based on contraction equations are applied in \cite{sato_linear_2017,tsilionis_tensor-based_2024}, where recursive linear equations are derived.